% ================================================================
%  第 5 章 劳动与发现 (原书 p.127-132, 边码 p.1-p.8)
%  源文: _work/src1/task-ch05.md
% ================================================================
% ---------------- p.127 ----------------
\chapter{劳动与发现}

\noindent 1983年6月\par

\section{(1) 孩子与好上帝}

\origpage{1} 我眼下正在做的这些数学笔记,是十三年来头一批我打算拿去发表的文字。读者不会奇怪:沉默了这么久之后,我的表达风格变了。然而这种表达上的变化,并不是工作风格或工作方法有什么变化的标志\fnoteP{1}{(1)}{(1984年3月补写)说我的工作“风格”与“方法”没有变,恐怕失于武断——毕竟我在数学上的表达风格已经深刻改变了。这一年来用于《追寻栈》的大部分时间,是坐在打字机前敲下一些思考;它们就是要几乎原样发表的(只除了后来补加的若干相当简短的注,借互见、勘误之类方便阅读)。没有剪刀,没有浆糊,不必费劲去炮制一份“定稿”(那份定稿尤其不该让任何通向它的过程透出影来)——这毕竟也是“风格”与“方法”的变化!除非把数学工作本身与写作、与陈述结果的工作割裂开来,而那是人为的,因为那不合事情的实情:数学工作与写作是拆不开的。} (1),更不是说我的数学工作本身的性质发生了转变。不但这性质一如既往——而且我确信:发现的工作,其性质在一个人与另一个人身上是相同的,它超越种种制约与气质所造成的差异,那些差异千变万化,无穷无尽。

发现是孩子的特权。我要说的是小孩子——那个还不怕弄错的孩子,不怕显得像个傻瓜,不怕不正经,不怕不跟大家一样。他也不怕他瞧着的东西不识趣:与他对它们的期待不一样、与它们本该是的模样不一样,或者更确切地说:与人们理所当然认定它们是的模样不一样。他不懂那些无声而毫无缺口的共识——它们是我们呼吸的空气的一部分,是所有明白人、所有公认的明白人的空气。天知道,自时间的黑夜深处以来,出过多少这样的人——明白人,公认的明白人!

我们的头脑里塞满了一种杂七杂八的“知识”——恐惧与懒惰、渴念与禁忌纠缠成一团;还有随处涌来的信息和按一按就出来的解释——一个封闭的空间,信息、渴念与恐惧在里面堆积,而旷野的风从来吹不进去。除了那点例行公事的“会做”,这“知识”的主要作用,似乎就是排掉活的知觉,排掉对世间事物的认识。它的效果,主要是一种巨大的惰性,一种往往压得人喘不过气来的重负。

% ---------------- p.128 ----------------
小孩子发现世界,就像他呼吸——呼吸的涨落让他把世界接进自己娇嫩的存在里,也让他投身于接纳他的世界。成年人也在发现,在那些稀少的时刻:他忘了自己的恐惧和知识,用睁得大大的、渴于认知的眼睛看事物、看自己,一双崭新的眼睛——一双孩子的眼睛。

\asterisms

\origpage{2} 上帝一边发现世界,一边创造世界;更确切地说,他永恒地\textbf{创造}着世界,一边创造一边发现——也一边发现一边创造。他创造了世界,又日复一日地创造着它,反反复复,千百万个千百万次,不歇息,摸索着,错了千百万个千百万次又拨正航向,不知疲倦……每一次,在这场用探针敲击万物、听万物回应的游戏里(它们回答:“这一下敲得不赖”,或者:“这一下你可是彻底犯浑”,或者:“顺顺当当,就这么来”),在又一次探针中——它修正或重提前一次探针,回应前一个回答……在造物主与万物之间这场无限对话的每一次往复里(这对话在创世的每一刻、每一处发生),上帝学习着、发现着,他越来越亲切地认识万物,而万物也在他手中渐渐有了生命、有了形状、不断变形。

这就是发现与创造的路数;亘古以来,似乎一直如此(就我们所能认识的而言)。一直如此——并不需要人类姗姗来迟地登场(不过是一两百万年前的事),撸起袖子插手进来——再加上最近那些众所周知的糟糕后果。

我们当中,不时有人发现这样或那样的事物。有时他会在自己的生命里重新发现——怀着惊叹——\textbf{发现}究竟是怎么一回事。每个人身上都有他所需要的一切,去发现这广袤世界上吸引他的一切,包括他身上那份奇妙的能力——世上最简单、最显然不过的东西!(然而许多人已经忘了的东西,就像我们忘了唱歌,忘了像孩子那样呼吸……)

每个人都能重新发现发现与创造是怎么一回事,而没有人能把它发明出来。它们先于我们就在那里,它们就是它们之所是。

\section{(2) 错误与发现}

回到我真正意义上的数学工作的“风格”,或它的“性质”、它的“路数”:它们如今与从前一样,都是好上帝亲自教给我们的——不用言语,教给每一个人,天知道在什么时候,也许远在我们出生之前很久。我照他的样子做。这也是每个人出于本能所做的,一旦好奇心驱使他去认识某一样东西——从那一刻起,这东西就被这份欲望、这份渴念占据了……

\origpage{3} 当我对一样东西感到好奇,数学的也好,别的也好,我就去问它。我去问它,不在意我的问题也许很蠢、或者会显得很蠢,也不要求它非得先经过反复掂量。问题常常采取断言的形式——一个其实只是探针的断言。我对自己的断言信几分?这当然取决于我在理解眼前之物上走到了哪一步。往往,尤其在一项研究的开头,这断言干脆就是错的——可还是得把它做出来,才能让自己信服它是错的。往往,只消把它写下来,
% ---------------- p.129 ----------------
错就一目了然;而在写下来之前,只有一片模糊,一种不适,而不是这份明白。这样就能少一分无知,再次进攻,带着一个也许不那么“离谱”的问题—断言。更常见的是:断言照字面理解是错的,但那个仍旧笨拙地试图借它表达出来的直觉是对的,虽然还模糊。这直觉会渐渐从一层同样不成形的脉石中析出——起初是些错误或不恰当的观念;它会渐渐走出“未被理解”的幽冥(那只是等着被理解),走出“未知”的幽冥(那只是等着被认识),取得一个只属于它自己的形式,让轮廓变得精细、鲜明——随着我向眼前这些事物提出的问题越来越精确、越来越切题,把它圈得越来越近。

但也会这样:靠这条路子,反复的探针会聚成局势的某一幅图像,从雾中显出相当清晰的线条,足以让人开始相信这幅图像很好地表达了实情——其实满不是那么回事,只要这图像沾上一个分量十足的错,一个足以叫它深深走样的错。从最初发现的“对不上号”出发——所得图像与某些确凿事实对不上,或与另一些同样为我们所信任的图像对不上——逐步查出这样一个错误观念的工作,有时很艰苦,而且往往伴随着越来越紧的张力,随着人逼近矛盾的结,那矛盾先还模糊,渐渐越来越刺耳——直到它终于爆开:错误的发现,某一种世界观的坍塌,像一次解放。\origpage{4}\textbf{错误的发现,是一切发现工作中最关键的时刻之一,是其中最富创造性的时刻之一},无论那是数学工作,还是发现自我的工作。正是在那一刻,我们对所探究之物的认识突然更新了。

\textbf{怕错与怕真,是同一回事。}害怕弄错的人,无力发现。正是在我们害怕弄错的时候,我们身上的错误才变得岩石一般不可动摇。因为在恐惧中,我们死死抓住某一天被我们判定为“真”的东西,或者抓住那些历来就被当作“真”呈现给我们的东西。当我们不是被害怕失去一种虚幻安全的恐惧所驱动,而是被求知的渴念所驱动,那时错误就像痛苦或悲伤一样,从我们身上穿过,永不定形,而它经过的痕迹,是一种更新了的认识。

\section{(3) 不可告人的劳作}

一切真正研究的自发路数,可以说从来不出现在那些本应传达、传递“被找到”之物的实质的文本或话语里——这绝不是偶然。文本与话语,多半只满足于记下一些“结果”,用的形式,在凡夫俗子眼里,必须显得像一条条严峻而不变的定律,亘古以来就刻在某种巨型图书馆的花岗岩法版上,由某位全知的神口授给那些入门者—书记—学者之流;口授给那些写学术书和同样学术的论文的人,那些站在讲坛高处、或在讨论班更小的圈子里传递知识的人。有没有哪怕一本教科书、一本供小学生、中学生、大学生乃至“我们的研究者”用的手册,能给不幸的读者一点点关于研究究竟是什么的概念——除了那个普遍接受的概念:研究嘛,就是当你很有学问、考过一大堆考试甚至竞考的时候,就是那些大脑袋,巴斯德和居里,还有诺贝尔奖得主什么的……而我们这些读者或听众,勉勉强强把大人物们为了人类的利益而惠予记下的“知识”吞下去,我们顶多(如果肯下苦功)够格通过年终考试,而且
% ---------------- p.130 ----------------
还不一定……

\origpage{5} 这样的人有多少呢——包括那些不幸的“研究者”自己,那些为论文或文章发愁的人,甚至包括我们当中最有“学问”、最有声望的人——有谁天真到看得出:“寻找”不是别的,正是怀着激情\textbf{去问}万物——就像一个孩子\textbf{想知道}他自己或他的小妹妹是怎么来到这世上的。寻找与找到,也就是:提问与倾听,是世上最简单、最自然不过的事,世上没有人拥有它的特权。这是我们一出生就全都收到的一份“馈赠”——它注定要以无穷无尽的面孔表达自己、舒展自己,此一时与彼一时不同,一个人与另一个人也不同。

当你冒险把这类话说出口,你在这一位和那一位那里收到的,从最确信自己是笨学生的笨学生,到最确信自己有学问、远在凡夫之上的大学者,都是一样的笑:一半尴尬,一半心照不宣,好像你刚讲了一个有点粗俗的笑话,好像你在显摆一种一眼就能看穿的天真;“这话都不错,谁也不该看不起,这我同意——可也别太过分啊——笨学生就是笨学生,他不是爱因斯坦,也不是毕加索!”

面对这样一致的赞同,我再坚持就不识趣了。真是不可救药,我又一次失去了闭嘴的机会……

不,这绝不是偶然:各类教育书、劝世书和各种手册,异口同声地把“知识”呈现得像是一出娘胎就穿得整整齐齐,从那些天才的头脑里出来——那些头脑为了我们的益处把它记了下来。也不能说这是恶意,即便在少数场合,作者相当“门儿清”,知道那幅图像(他的文本免不了要暗示的图像)与实情毫不相符。在这种场合,陈述有时不止是一堆结果和诀窍:有一股气息穿过它,有一种活的视野赋予它生命,而这视野有时便从作者传到专心的读者那里。但一种默契——力量看来相当可观——使得文本不留下半点它由之产生的那些工作的痕迹,即使它以镌刻般的笔力表达出对事物有时相当深刻的视野,而那视野正是这工作真正的果实之一。

\origpage{6} 说实话,有些时候,我自己也隐隐感到这股力量、这份无声共识的分量——就在我筹划写作、出版《数学反思》的时候。如果我试着去探一探这份共识所采取的默不作声的形式,或者更确切地说,去探一探我内心对我这计划的抵抗(它由这份共识引发)所采取的形式,立刻冒出来的词是“不体面”。这份共识,不知从何时起已内化在我心里,它对我说(而这是我头一回费心把它这几周、也许是这几个月来带几分固执地对我嘟囔的东西,拉到日光下,拉进我的视野里):“把一项发现工作中的起起落落、边角上那些弄砸了的摸索、总之那些‘脏衣服’,摊在别人面前,甚至当众摊开——这不体面。这只是浪费读者的时间,而读者的时间是宝贵的。再说,这还要多出好多页,得排版、得印刷——多浪费啊,科学印刷品价钱这么贵!非得虚荣到什么地步,才会这样摊开那些对谁都没意思的东西,好像连我那些胡乱的折腾也成了了不起的事——总而言之,一个摆谱的机会。”还有更隐秘的一层:“发表这样一份反思的笔记——照它真正进行的样子——是不体面的,就像在公共广场上做爱是不体面的,或者把分娩的劳作留下的血污床单拿出来展示、哪怕只是让它晾在那里,也是不体面的……”

这里的禁忌,采取的是性禁忌那副既阴险又专横的面孔。正是在写这篇导论的时候,我才开始隐约看见它非凡的力量,以及这本身即非凡的事实所含的意义——这事实正是那股力量的证明:发现的真正路数,如此
% ---------------- p.131 ----------------
令人不知所措地简单,一种孩子气的简单,却几乎无处透出;它被无声地抹去、忽视、否认。甚至在科学发现这个相对无害的领域里也是如此——谢天谢地,那不是发现自己的“小鸡鸡”之类的事——总之是一项可以放心交到所有人手里的“发现”,而且(人们会以为)没什么可隐瞒的……

倘若我愿意顺着这里出现的这根“线”走下去——这根线一点也不细,反而粗壮结实得很——它肯定会把我带到远得多的地方,远过那几百页同调—同伦代数;那些页面我终究会写完,交给印刷厂的。

\section{(4) 他人之无谬与对己之轻蔑}

\origpage{7} 刚才我谨慎地说“我的表达风格”变了,还让人以为这没什么可奇怪的——说到底,那是个委婉说法。您明白:十三年没写过东西,当然跟从前不一样了,“表达风格”总得变,必然的……区别在于,从前我“表达”(原文如此)得像所有人一样:我先做工作,然后倒过来重做一遍,把所有的涂改都仔细擦掉。擦着擦着又出现新的涂改,把整个工作搅得有时比初稿还糟。于是又得重做——有时三遍,甚至四遍,直到一切都无可挑剔。不但没有可疑的角落,也没有悄悄扫到合宜的家具底下的垃圾(我从来就不喜欢角落里的垃圾,既然都费心扫地了);尤其是,读定稿的时候,它给人的印象固然体面(像任何别的科学文本一样),那就是:作者(此处即区区在下)是化为肉身的无谬。他万无一失地正好落在“那些”正确的概念上,然后是“那些”正确的陈述上,一环扣一环,像上了油的马达一样低吟着,那些证明“咚”的一声落下来,每一个都恰好在它该来的时候!

\origpage{8} 效果如何,请设想一个毫无戒心的读者:比如说一个正在学毕达哥拉斯定理或二次方程的中学生,甚至我的一位在所谓“高等”研究或教学机构里的同行(明白人一听就懂,此致!),(不妨说)正费劲地读某位显赫同行的某篇文章!这类经验在一个小学生——甚至大学生或研究者——的一生中重复成百上千次,再加上家庭里以及全世界所有国家的所有媒体里那场合宜的合唱,其效果可想而知。只要肯留心,人在自己身上、在别人身上都能看到:那是一种对自身毫无价值的内心确信,与那些“懂行的人”、“干事的人”的能力和重要性恰成对照。

这份内心的确信,有时会被一种能力所补偿——记住不理解之物的能力,甚至某种操作上的灵巧:矩阵相乘,用“正题”和“反题”拼出一篇法文作文……总之是鹦鹉或耍把戏的猴子的本领,如今比以往任何时候都更受看重,由人人艳羡的文凭认可,由安稳的职业生涯报偿。但即便那位文凭缝满一身、位置坐得安稳、也许还满身荣誉的人,在心底里也并不上当,并不把这些重要、这些“价值”的假标志当真。甚至那位更少见的人——他把全部赌注押在某样真正天赋的发展上,在职业生涯中懂得施展身手、做出创造性作品——在心底里也并不信服自己名声的光彩,尽管他常常想借这光彩让自己、也让别人看走眼。一种从未被审视过的同样的怀疑,住在这一位和那一位身上,也住在随便哪个笨学生身上;一种同样的确信,也许他们永远不敢去认识它。

正是这份怀疑、这份未曾说出的内心确信,推动着这一位和那一位不断地超越自己——
% ---------------- p.132 ----------------
在荣誉或作品的积累中;并且把他们私下里啃噬着自己的那份对己之轻蔑,投射到别人身上(尤其是那些他们对其多少有些权力的人身上……)——这是一场注定不可能的逃脱尝试,靠的是积累他们优于别人的“证据”\fnoteP{2}{(2)}{(1984年3月补写)重读最后这两段,我有些不适之感,因为写它们的时候,我牵连的是别人,而不是我自己。显然,写的时候,“我自己也可能与此有关”这个念头并没有碰到我。我肯定什么也没学到:我就这样满足于把多年来在别人身上察知、并以许多方式看着被证实的东西,白纸黑字写下来(无疑还带着某种满意)。在这份反思的后续部分,我不得不回想起:我的一生中,对别人的轻蔑姿态并不缺少。我所把握到的他人之轻蔑与对己之轻蔑之间的关联,若在我这个人身上竟然缺席,那才奇怪;健全的理性(以及类似情形下我对自身视而不见的经验——我终究还是察觉到了)告诉我:肯定不该是这样!然而眼下这只是一条纯粹的推论,它唯一可能的用处,是促使我亲眼看一看究竟发生了什么,去看、去审视(如果它确实存在过、或还存在着)这份仍属假设的对己之轻蔑——它埋得如此之深,至今完全逃过了我的目光。真的,可看的东西并不缺!这一件忽然在我看来是最要紧的几件事之一,恰恰因为它藏得如此之深。……\emph{[(1984年8月)不过关于这一点,参见“屠杀”一注(n\textsuperscript{o}~87)末两段的反思。]}} (2)。
